\section{Planteamiento del problema}
\label{sec:problema}

La búsqueda de empleo remoto es hoy una experiencia fragmentada. Las ofertas se dispersan entre decenas de plataformas —agregadores generalistas, portales especializados y bolsas de empleo corporativas—, cada una con su propio esquema de filtros y su propio motor de búsqueda. La herramienta subyacente que casi todas comparten, sin embargo, ha cambiado poco en las últimas dos décadas: el emparejamiento por palabra clave (\textit{keyword-matching}) entre los términos de la consulta y los campos estructurados o el texto libre de la oferta.

Esta dependencia del emparejamiento léxico exacto es también la base de la mayoría de los sistemas de seguimiento de candidatos (\textit{Applicant Tracking Systems}, ATS) usados del lado del reclutador. La propia motivación de trabajos como el de \citet{HaThuc2017} —que documenta la migración de LinkedIn desde una búsqueda de talento basada en palabras clave (\textit{Query-By-Keyword}) hacia una basada en ejemplos (\textit{Query-By-Example})— y el de \citet{Bocharova2023} —que atribuye explícitamente a la variación léxica entre títulos de vacantes equivalentes ("Software Engineer" frente a "Backend Developer", por ejemplo) una pérdida sistemática de coincidencias relevantes— evidencia que el problema no es anecdótico sino estructural: el emparejamiento exacto falla precisamente en los casos en que existe sinonimia entre habilidades, variación en los títulos de cargo o ambigüedad en la formulación de la consulta. La revisión bibliométrica de \citet{RojasGaleano2022} sobre investigación en inteligencia artificial aplicada al emparejamiento empleo-currículo confirma que esta limitación ha sido, durante más de una década, uno de los motivadores centrales de la investigación en la intersección de la recuperación de información y el reclutamiento asistido por cómputo.

Frente a este panorama, la generación aumentada por recuperación (\textit{Retrieval-Augmented Generation}, RAG), propuesta originalmente por \citet{Lewis2020}, ofrece una alternativa arquitectónica distinta: en lugar de exigir una coincidencia léxica exacta, el sistema recupera los fragmentos de información semánticamente más cercanos a la consulta —representados como vectores en un espacio de \textit{embeddings}— y delega en un modelo de lenguaje la síntesis de una respuesta fundamentada en dichos fragmentos. Como documenta la revisión de \citet{Gao2023}, este paradigma ha probado ser particularmente efectivo para reducir la alucinación de los modelos de lenguaje y para responder preguntas que dependen de una base de conocimiento externa y actualizable, dos propiedades directamente aplicables al dominio de la búsqueda de empleo: el conocimiento (las ofertas vigentes) cambia constantemente, y una respuesta no fundamentada en datos reales (por ejemplo, un salario o un requisito inventado) es inaceptable para quien la usa para tomar una decisión de carrera.

La mayoría de las plataformas de empleo que hoy incorporan "inteligencia artificial" lo hacen de forma superficial: un asistente conversacional genérico superpuesto a los mismos filtros de siempre, o un puntaje de similitud aislado que no reemplaza al motor de búsqueda. JobAI —el artefacto sobre el que se desarrolla este trabajo (véase el Capítulo 4)— se construyó bajo una premisa distinta: que el copiloto de IA \emph{sea} el motor de búsqueda, mediante un pipeline RAG completo que recupera, filtra y sintetiza información sobre ofertas laborales agregadas de múltiples fuentes, en lugar de limitarse a conversar sobre datos que el usuario ya podría ver en un panel de filtros. Sin embargo, aunque JobAI ya existe como artefacto funcional, dos vacíos persisten y motivan el presente trabajo. Primero, un vacío de evaluación: no existe todavía una medición cuantitativa, sistemática, de qué tan bien cumple el pipeline RAG su promesa de precisión y fidelidad frente al enfoque de búsqueda por palabra clave que busca superar. Segundo, un vacío de fundamentación arquitectónica: la decisión de almacenamiento vectorial —inicialmente resuelta con la extensión pgvector sobre PostgreSQL por su simplicidad operativa— no ha sido contrastada empíricamente frente a una base de datos vectorial dedicada, a pesar de que la literatura reciente en sistemas de gestión de bases de datos vectoriales (revisada en el Capítulo 2) documenta de manera consistente las limitaciones de escalabilidad de la primera aproximación.

\section{Formulación del problema}

De lo anterior se deriva la pregunta de investigación que orienta este trabajo:

\begin{quote}
\textit{¿En qué medida un pipeline de Retrieval-Augmented Generation, implementado sobre una plataforma real de búsqueda de empleo remoto y evaluado mediante métricas cuantitativas de precisión de recuperación, fidelidad de las respuestas y latencia, permite superar las limitaciones del emparejamiento por palabra clave, y cómo debe diseñarse dicha evaluación para que sus resultados sean válidos y reproducibles?}
\end{quote}

Esta pregunta tiene una segunda dimensión, subordinada pero necesaria para responderla con rigor metodológico: dado que JobAI ya es un artefacto construido, ¿cómo se enmarca metodológicamente la evaluación de un artefacto preexistente, y qué papel juega dentro de ese marco la decisión de migrar el almacenamiento vectorial de pgvector a una base de datos vectorial dedicada (Qdrant)? El Capítulo 3 retoma esta segunda dimensión en detalle, apoyándose en el marco de Design Science Research introducido conceptualmente en la Sección~\ref{sec:marco-teorico-dsr} de este capítulo y revisado críticamente en el Capítulo 2.

\section{Objetivos}

\subsection{Objetivo general}

Desarrollar y evaluar una plataforma de búsqueda de empleo remoto que utiliza un pipeline de Retrieval-Augmented Generation para permitir búsqueda semántica y conversacional sobre ofertas laborales, superando las limitaciones de precisión de la búsqueda por palabras clave.

\subsection{Objetivos específicos}

\begin{enumerate}
    \item Caracterizar, a partir de la literatura, las limitaciones de los sistemas actuales de búsqueda de empleo y las oportunidades que ofrece RAG para superarlas.
    \item Diseñar y documentar la arquitectura del pipeline RAG (ingestión, \textit{chunking}, \textit{embeddings}, clasificación de consultas, síntesis) aplicada al dominio de búsqueda de empleo, incluyendo la migración del almacenamiento vectorial de pgvector a una base de datos vectorial dedicada.
    \item Definir una metodología de evaluación con métricas cuantitativas (precisión de recuperación, fidelidad/\textit{groundedness} de las respuestas, tiempo de respuesta) para validar la plataforma.
    \item Evaluar los resultados obtenidos y contrastarlos frente a un enfoque de búsqueda tradicional por palabras clave.
\end{enumerate}

\section{Justificación}

La justificación de este trabajo se sostiene sobre tres ejes complementarios.

Desde el punto de vista \textbf{práctico}, la fragmentación y las limitaciones de precisión descritas en la Sección~\ref{sec:problema} tienen un costo directo para quien busca empleo: tiempo perdido evaluando ofertas irrelevantes y oportunidades genuinamente compatibles que nunca aparecen porque el sistema de búsqueda no reconoció la equivalencia semántica entre lo que la persona escribió y lo que la oferta describe. Una plataforma que resuelve esto con evidencia cuantitativa de su desempeño —y no solo con la promesa genérica de "búsqueda con IA"— tiene valor inmediato y verificable.

Desde el punto de vista \textbf{teórico}, la revisión bibliométrica de \citet{RojasGaleano2022} muestra que la investigación en inteligencia artificial aplicada al emparejamiento empleo-currículo es un campo activo pero todavía concentrado en tareas de \textit{ranking} y recomendación (como los sistemas de \citealp{HaThuc2017,Shen2024,Bocharova2023} revisados en el Capítulo 2); la aplicación específica de arquitecturas RAG completas —con recuperación, síntesis conversacional y clasificación híbrida de consultas— al dominio de búsqueda de empleo es todavía un espacio de trabajo emergente. Este trabajo aporta, en ese sentido, un caso de estudio evaluado empíricamente en un subdominio donde la literatura de RAG general \citep{Lewis2020,Gao2023} y la literatura específica de matching de empleo rara vez se han integrado explícitamente.

Desde el punto de vista \textbf{metodológico}, este trabajo constituye una aplicación del paradigma de Design Science Research \citep{Hevner2004,Peffers2007} a un artefacto de software ya construido y en producción, en lugar de a un artefacto diseñado desde cero para fines exclusivamente académicos. Esto exige adaptar el ciclo de diseño-evaluación de Peffers et al. a un escenario de iteración sobre un sistema existente —la migración del almacenamiento vectorial y el diseño de la evaluación cuantitativa son, en los términos de ese marco, una nueva iteración del ciclo—, lo cual se documenta en detalle en el Capítulo 3 y puede sensiblemente extenderse a un artefacto de similar naturaleza en trabajos de investigación aplicada.

\section{Alcance y limitaciones}

El alcance de este trabajo se circunscribe a JobAI como artefacto de estudio: una plataforma que agrega ofertas de empleo remoto de tres fuentes públicas (Remotive, Jobicy y Findwork), las indexa mediante \textit{embeddings} vectoriales y expone un copiloto conversacional basado en un pipeline RAG (véase el Capítulo 4 para el detalle arquitectónico). La evaluación cuantitativa (Capítulo 5) se limita a las métricas de precisión de recuperación, fidelidad de las respuestas y latencia definidas en el Capítulo 3, contrastadas contra un \textit{baseline} de búsqueda por palabras clave sobre el mismo conjunto de datos.

Quedan fuera del alcance: (i) una evaluación de experiencia de usuario o de adopción con usuarios reales, dado que el enfoque es la calidad técnica del pipeline y no su recepción; (ii) la generalización de los resultados a otros dominios de búsqueda distintos al de empleo remoto; (iii) la comparación exhaustiva contra la totalidad de bases de datos vectoriales disponibles en el mercado, restringiéndose la comparación empírica a pgvector y Qdrant como representativas de los dos enfoques arquitectónicos discutidos en el Capítulo 2 (extensión vectorial sobre un motor relacional frente a motor vectorial dedicado); y (iv) la evaluación de la calidad de los datos de origen (Remotive, Jobicy, Findwork), que se asumen como dados.

\section{Marco teórico conceptual}

Esta sección introduce, a nivel conceptual, el vocabulario técnico necesario para seguir el resto del documento. La revisión crítica de cada una de estas líneas —con sus fuentes primarias, sus hallazgos específicos y su relación explícita con las decisiones de diseño de JobAI— se desarrolla en el Capítulo 2.

\subsection{Recuperación de información y el problema del emparejamiento léxico}

La recuperación de información clásica enfrenta un problema bien documentado: el \textit{desajuste de vocabulario} (\textit{vocabulary mismatch}), en el que dos textos semánticamente equivalentes no comparten los términos exactos necesarios para que un sistema de coincidencia léxica los relacione. En el dominio de reclutamiento, este problema se manifiesta como sinonimia entre habilidades (p. ej., "JS" y "JavaScript"), variación entre títulos de cargo equivalentes, y ambigüedad en la forma en que una persona describe lo que busca frente a cómo una oferta describe lo que ofrece. Los sistemas de búsqueda por palabra clave y los ATS tradicionales, al depender de la coincidencia exacta (o de expansión limitada mediante sinónimos predefinidos), heredan esta limitación de forma directa.

\subsection{Modelos de lenguaje de gran escala}

Los modelos de lenguaje de gran escala (\textit{Large Language Models}, LLMs) son modelos generativos preentrenados sobre grandes volúmenes de texto, capaces de producir lenguaje natural fluido y de razonar sobre instrucciones formuladas en lenguaje natural. Su limitación central, señalada de forma recurrente en la literatura de RAG \citep{Gao2023}, es que su conocimiento está codificado paramétricamente en los pesos del modelo en el momento del entrenamiento: no tienen acceso nativo a información posterior a su corte de entrenamiento ni a datos privados o específicos de un dominio (como el catálogo de ofertas de una plataforma de empleo), y tienden a "alucinar" —generar afirmaciones plausibles pero no sustentadas— cuando se les pide responder sobre información que no poseen.

\subsection{Retrieval-Augmented Generation (RAG)}

RAG, propuesto por \citet{Lewis2020}, combina un componente de \textit{recuperación} (que busca los fragmentos de una base de conocimiento externa más relevantes para una consulta dada, típicamente mediante similitud vectorial) con un componente \textit{generativo} (un LLM que sintetiza una respuesta a partir de la consulta y los fragmentos recuperados). La ventaja central de este diseño es que ancla la respuesta del modelo en información verificable y actualizable, en lugar de depender exclusivamente de su conocimiento paramétrico. La encuesta de \citet{Gao2023} organiza la evolución de este paradigma en tres niveles de sofisticación creciente —RAG ingenuo (recuperación y generación en un solo paso), RAG avanzado (con reformulación de consultas, reordenamiento de resultados y filtrado de contexto) y RAG modular (con componentes intercambiables y flujos condicionales, como enrutamiento entre distintas estrategias de recuperación)—, una taxonomía que resulta directamente pertinente para ubicar el diseño del pipeline de JobAI (Capítulo 4). Evaluar un sistema RAG plantea, a su vez, un reto propio: al no existir siempre una respuesta de referencia contra la cual comparar, \citet{Es2024} proponen con RAGAS un conjunto de métricas de evaluación sin referencia humana —fidelidad (\textit{faithfulness}, si la respuesta se sostiene en los fragmentos recuperados), relevancia de la respuesta, y precisión/exhaustividad del contexto recuperado— que este trabajo retoma como base conceptual para el diseño de las métricas de evaluación del Capítulo 3.

\subsection{Bases de datos vectoriales y búsqueda por similitud}

Un \textit{embedding} es una representación vectorial de un fragmento de texto (o de cualquier otro dato) en un espacio de alta dimensión, construida de forma que la distancia (típicamente similitud coseno) entre dos vectores refleje la similitud semántica entre los textos que representan. Buscar los fragmentos más relevantes para una consulta se convierte, en este espacio, en un problema de búsqueda del vecino más cercano (\textit{nearest neighbor search}). Dado que la búsqueda exacta de vecinos más cercanos es computacionalmente costosa a gran escala, en la práctica se recurre a algoritmos de búsqueda \textit{aproximada} del vecino más cercano (ANN); el más influyente de ellos es HNSW \citep{Malkov2020}, una estructura de grafo jerárquico de navegación que permite búsquedas en tiempo logarítmico. Estos algoritmos pueden implementarse tanto como una extensión sobre un motor de base de datos relacional (como pgvector sobre PostgreSQL) como dentro de un motor construido específicamente para datos vectoriales (como Qdrant o Milvus). El Capítulo 2 revisa en detalle la evidencia empírica que sustenta la decisión de migrar de la primera aproximación a la segunda.

\subsection{Design Science Research como paradigma metodológico}
\label{sec:marco-teorico-dsr}

Design Science Research (DSR) es un paradigma de investigación en sistemas de información orientado a la creación y evaluación rigurosa de artefactos —métodos, modelos, constructos o instanciaciones de software— que resuelven problemas organizacionales o técnicos relevantes. \citet{Hevner2004} formalizan este paradigma mediante siete lineamientos (entre ellos, que el artefacto debe ser evaluado rigurosamente y que la investigación debe producir una contribución verificable), mientras que \citet{Peffers2007} lo operacionalizan en un proceso de seis actividades secuenciales —identificación del problema, definición de objetivos, diseño y desarrollo, demostración, evaluación y comunicación— que admite entradas en distintos puntos del ciclo, incluyendo la iteración sobre un artefacto ya existente. Este trabajo adopta DSR como marco metodológico general; su aplicación concreta a JobAI —incluyendo cómo se interpreta la migración a Qdrant como una nueva iteración del ciclo diseño-evaluación— se desarrolla en el Capítulo 3.
